% Contexto e instrucciones literales de las páginas 1 y 2 de EXA-0054.
% Fuente: Selectividad/sel_2011_matematicas_aplicadas/121-EXAMEN 4.pdf
% SHA-256: c70ed065da193c82dd0a02fb4c05069415b0adf53d0f49bba1b0375a19823566
% No se reproduce el logotipo ni el marco de la cabecera.
\section*{Contexto del examen original}
\textbf{UNIVERSIDADES DE ANDALUCÍA}\par
\textbf{PRUEBA DE ACCESO A LA UNIVERSIDAD}\par
CURSO 2010-2011\par
\textbf{MATEMÁTICAS APLICADAS A LAS CIENCIAS SOCIALES II}

\textbf{Instrucciones:}
\begin{enumerate}[label=\textbf{\alph*)}]
  \item Duración: 1 hora y 30 minutos.
  \item Elija una de las dos opciones propuestas y conteste los ejercicios de la opción elegida.
  \item En cada ejercicio, parte o apartado se indica la puntuación máxima que le corresponde.
  \item Se permitirá el uso de calculadoras que no sean programables, gráficas ni con capacidad para almacenar o transmitir datos.
  \item Si obtiene resultados directamente con la calculadora, explique con detalle los pasos necesarios para su obtención sin su ayuda. Justifique las respuestas.
\end{enumerate}
\noindent\textbf{Nota de la colección:} Se reúnen ambas opciones para consulta. La regla anterior pertenece al examen original. Curso 2010-2011 impreso. Convocatoria,tipo,reserva y modelo no confirmados;121 y4 sólo del nombre. Examen y criterios sin alias exactos. Tabla normal SRC-0DF80B854632 asociada por prefijo121 y carpeta;cuatro procedencias idénticas. INC-020 conserva la errata literal de A4.
\par\medskip
